\chapter{Solutions}

\runinhead{Démonstration}\phantomsection\addcontentsline{toc}{subsubsection}{Démonstration}\label{IV-answer-bienaymuxe9cheby}
C'est une application immédiate de
\protect\hyperref[IV-inegmarkov]{l'inégalité de
Markov (p.~\pageref*{IV-inegmarkov})} à \((X-\mathbb{E}(X))\) avec
\(p =2\).

\runinhead{Pile ou face
infini}\phantomsection\addcontentsline{toc}{subsubsection}{Pile ou face
infini}\label{IV-answer-pf}
En suivant l'indication, on a \(\mathbb{P}(A_n) = p\), d'où
\(\sum_{n=1}^\infty \mathbb{P}(A_n) = \infty\). Le lemme de
Borel-Cantelli nous indique alors que
\(\mathbb{P}(\lim \sup_{n \to \infty} A_n) = \mathbb{P}\left(\cap_{n\geq 1} \cup_{k \geq n} A_n \right) = 1\).
Autrement dit, on a presque sûrement un nombre infini de faces. En
passant au complémentaire, on en déduit que l'événement \(A\) est de
probabilité nulle.

\runinhead{En proba et en
moyenne}\phantomsection\addcontentsline{toc}{subsubsection}{En proba et
en moyenne}\label{IV-answer-lien}
Pour tout \(\varepsilon \in ]0,1[\), la probabilité
\(\mathbb{P}(|X_n|\geq \varepsilon) = \frac{1}{n}\) tend vers 0 quand
\(n\) tend vers l'infini. Ainsi, la suite
\((X_n)_{n\in \mathbb{N}^\ast}\) tend vers \(X=0\) en probabilité. Comme
\(\mathbb{E}(X_n) = \frac{1}{n}\), elle tend également en moyenne vers
0.

\runinhead{En proba mais pas en
moyenne}\phantomsection\addcontentsline{toc}{subsubsection}{En proba
mais pas en moyenne}\label{IV-answer-lien2}
Par le même argument que ci-dessus, nous voyons que la suite
\((Y_n)_{n\in \mathbb{N}^\ast}\) converge en probabilité vers 0, mais
comme \(\mathbb{E}(Y_n) = n\), la suite ne converge pas en moyenne vers
0 (ni vers aucune autre limite finie).

\runinhead{Presque
sûre}\phantomsection\addcontentsline{toc}{subsubsection}{Presque
sûre}\label{IV-answer-lien3}
Si \(\omega \in \{U > 0\}\) est fixé, alors
\(\exists n_0 \in \mathbb{N}\) tel que \(U(\omega) > \frac{1}{n_0}\), et
donc tel que \(Z_n(\omega) = 0\) pour tout \(n \geq n_0\). Comme
\(\mathbb{P}(U>0)=1\), ceci montre que la suite
\((Z_n)_{n\in \mathbb{N}^\ast}\) converge presque sûrement vers 0.

\runinhead{Intervalles de longueur aléatoire
(1)}\phantomsection\addcontentsline{toc}{subsubsection}{Intervalles de
longueur aléatoire (1)}\label{IV-answer-ia}
\(\mathbb{E}(X_n) = \frac{1}{n}\), et la suite \(X_n\) tend vers
\(X = 0\) en moyenne, et donc en probabilité.

\runinhead{Intervalles de longueur aléatoire
(2)}\phantomsection\addcontentsline{toc}{subsubsection}{Intervalles de
longueur aléatoire (2)}\label{IV-answer-ia2}
Supposons que les \(A_n\) soient placés bout-à-bout, en recommençant en
0 chaque fois qu'on arrive au point 1. Il est clair que l'on parcourt
indéfiniment l'intervalle {[}0, 1{]} (car la série de terme général
\(1/n\) diverge). Ainsi la suite numérique \(X_n (\omega)\) ne converge
pour aucun \(\omega\), et on n'a pas \(X_n \to X\) presque-sûrement.
Cependant comme la série \(\sum_n 1/n^2\) converge, il s'en suit que
\(X_{n^2} \to X = 0\) presque-sûrement. En effet, posant
\(B_n =\{\omega \in \Omega, X_{n^2}(\omega)=1\}\), on a
\(\mathbb{P}(B_n) = \frac{1}{n^2}\) et
\(\sum_{n=1}^{+\infty} \mathbb{P}(B_n) = \frac{\pi^2}{6}\). Le
\protect\hyperref[IV-BC]{lemme de Borel Cantelli (p.~\pageref*{IV-BC})}
nous indique alors que \(\mathbb{P}(B) = 0\) où
\(B = \cap_{n\geq 1}\cup_{k\geq n} B_n\) et donc
\(X_{n^2}(\omega) \to 0,\, \forall \omega \notin B\).

\runinhead{Loi faible des grands
nombres}\phantomsection\addcontentsline{toc}{subsubsection}{Loi faible
des grands nombres}\label{IV-answer-wlln}
Par l'inégalité de Bienaymé-Chebyshev, on a
\[\mathbb{P}(|M_n-m| > \epsilon) \leq \frac{\mathbb{V}(M_n)}{\epsilon} = \frac{\sigma^2}{n\epsilon},\]
où \(\sigma^2 = \mathbb{V}(X_i), \forall i \in \mathbb{N}^\ast\).

\runinhead{Exponentielle}\phantomsection\addcontentsline{toc}{subsubsection}{Exponentielle}\label{IV-answer-expo}
On a
\(\log\left(\prod_{i=1}^n Y_i\right)^{1/n} = \frac{1}{n}\sum_{i=1}^n X_i \to m\)
p.s. par la \protect\hyperref[IV-lfgn]{loi forte des grands
nombre (p.~\pageref*{IV-lfgn})}. Puisque l'exponentielle est continue,
la \protect\hyperref[IV-propconv4]{proposition de
continuité (p.~\pageref*{IV-propconv4})} nous indique que
\(\left(\prod_{i=1}^n Y_i\right)^{1/n} \to e^m\) p.s.

\runinhead{Variance}\phantomsection\addcontentsline{toc}{subsubsection}{Variance}\label{IV-answer-variance}
Les variables aléatoires \(Y_i = (X_i - m)^2\) sont i.i.d. telles que
\(\mathbb{E}(Y_i)= \sigma^2\). Par la \protect\hyperref[IV-lfgn]{loi
forte des grands nombre (p.~\pageref*{IV-lfgn})}, on a
\[\lim_{n \to \infty} \frac{1}{n} \sum_{i=1}^n (X_i - m)^2 \to \sigma^2  \text{   p.s.}\]

\runinhead{Suite
gaussienne}\phantomsection\addcontentsline{toc}{subsubsection}{Suite
gaussienne}\label{IV-answer-suitegauss}
Soit \(f\) une fonction continue bornée sur \(\mathbb{R}\). On a
\begin{align*}
\mathbb{E}(f(X_n)) &= \int_\mathbb{R}f(y) \frac{1}{\sqrt{2\pi}\sigma_n} \exp\left(-\frac{y^2}{2\sigma^2_n}\right) dy \\
             &\xrightarrow[n \to \infty]{} \int_\mathbb{R}f(y) \frac{1}{\sqrt{2\pi}\sigma} \exp\left(-\frac{y^2}{2\sigma^2}\right) dy
\end{align*} où l'on a utilisé le théorème de convergence dominée
appliqué à la suite de fonctions
\(\frac{1}{\sqrt{2\pi}\sigma_n} \exp\left(-\frac{y^2}{2\sigma^2_n}\right)\).

\runinhead{Surbooking}\phantomsection\addcontentsline{toc}{subsubsection}{Surbooking}\label{IV-answer-surbooking}
Le nombre de passager \(S_{420}\) se présentant effectivement à
l'embarquement suit une loi binomiale \(\mathcal{B}(420,0.92)\). On veut
donc calculer
\(\mathbb{P}(S_{420} > 400) = 1 - \mathbb{P}(S_{420}<400)\). Le TCL nous
fournit l'approximation
\[\frac{S_{420} - np}{\sqrt{n p (1-p)}} = \frac{S_{420} - 386.4}{5.56} \approx Y \sim \mathcal{N}(0,1)\]
Ainsi, \(\mathbb{P}(S_{420}<400) \approx \Phi(2.44) \approx 0.99\) La
probabilité de surbooking est donc de 1\%.

\section{Inégalités de
concentration}\label{IV-inuxe9galituxe9s-de-concentration-1}

\runinhead{Question
1}\phantomsection\addcontentsline{toc}{subsubsection}{Question
1}\label{IV-answer-idc1}
L'inégalité de Chebyshev nous donne pour tout \(a>0\)
\[\mathbb{P}\left(\left|\frac{1}{n}\sum_{i=1}^n X_i - m \right| > a\right) \leq \frac{\mathbb{V}(X_1)}{n a^2}\]
Prenant, \(\frac{\mathbb{V}(X_1)}{n a^2} = \delta\), on obtient
\(a = \frac{1}{\sqrt{n\delta}}\), d'où le résultat. On retrouve au
passage la même vitesse de convergence que celle donnée par le TCL.

\runinhead{Question
2}\phantomsection\addcontentsline{toc}{subsubsection}{Question
2}\label{IV-answer-idc2}

\begin{enumerate}
\def\labelenumi{\roman{enumi})}
\tightlist
\item
  \(M(u) = \mathbb{E}(e^{u(X_1 -m)}) = \phi_{X_1}(-iu) = e^{u^2 \sigma^2/2}\),
  où \(\phi_{X_1}\) est la fonction caractéristique de \(X_1\).
\item
  On a \(e^{u(M_n-nm)} \geq e^{ua} 1_{S_n -nm \geq a}\), d'où
  \[\mathbb{P}(S_n -nm \geq a) \leq \frac{\mathbb{E}(e^{u(M_n-nm)})}{e^{ua}} = e^{-ua}M(u)^n\]
  par indépendance des \(X_i\).
\item
  \begin{align*}
  \mathbb{P}(|Y_n - m| \geq \varepsilon) &= \mathbb{P}(S_n - nm \geq n\varepsilon) + P(S_n - nm \leq -n\varepsilon) \\
                         &\leq e^{-nu\varepsilon} (M (u))^n + e^{-n(-u)(-\varepsilon)} (M (-u))^n = 2e^{-nu\varepsilon}(M (u))^n \\
                         &\leq 2e^{-nu\varepsilon}e^{nu^2\sigma^2/2}, \,\,\, \forall u \geq 0
  \end{align*} La meilleure majoration va être obtenue en minimisant
  l'exposant, c'est-à-dire pour \(u = \varepsilon/\sigma^2\). Nous en
  déduisons l'inégalité de Chernov.
\end{enumerate}

\runinhead{Question
3}\phantomsection\addcontentsline{toc}{subsubsection}{Question
3}\label{IV-answer-idc3}
Par l'inégalité de Bienaymé-Chebyshev,
\(\mathbb{P}(|Y_n - m| > \varepsilon) \leq \frac{\mathbb{V}(Y_n)}{\varepsilon^2} = \frac{\sigma^2}{n\varepsilon^2}\).

Ainsi, \(\mathbb{P}(|Y_n - m| \leq \varepsilon)\leq \alpha\) dès que
\(1-\frac{\sigma^2}{n\varepsilon^2}\). Avec \(\varepsilon = 0,05\), il
vient \(n \geq 80000\).

Par l'inégalité de Chernov, nous obtenons
\(2e^{- n\varepsilon^2/20} \leq 0, 05\). Il vient que \(n \geq 29512\).
Pour avoir une évaluation du même ordre de la probabilité cherchée, nous
pouvons donc prendre un échantillon beaucoup plus petit si nous
utilisons l'inégalité de Chernov. A taille d'échantillon fixée, nous
aurons une meilleure évaluation avec cette inégalité.

\section{Lemme de Borel-Cantelli}\label{IV-lemme-de-borel-cantelli-1}

\runinhead{Question
1}\phantomsection\addcontentsline{toc}{subsubsection}{Question
1}\label{IV-answer-bc1}
On voit dans un premier temps que
\(\bigcap_{n\geq 0} \bigcup_{k \geq n} A_k \in \mathcal{A}\) par unions
et intersections dénombrables. On a
\[\mathbb{P}(\lim \sup_n A_n ) = \lim_{p \to \infty} \mathbb{P}(\cup_{n\geq p} A_n) \leq \lim_{p \to \infty} \sum_{n \geq p} \mathbb{P}(A_n),\]
où on remarque que les deux suites sont décroissantes.

Si la série \(\sum_n \mathbb{P}(A_n)\) est convergente, le reste de
cette série tend vers 0 et l'inégalité implique que
\(\mathbb{P}(\lim \sup_n A n) = 0\).

\runinhead{Question
2}\phantomsection\addcontentsline{toc}{subsubsection}{Question
2}\label{IV-answer-bc2}
Supposons maintenant que les \(A_n\) soient indépendants et que la série
diverge. Soit \(m\) un nombre entier. Nous avons \begin{align*}
\mathbb{P}(\cup_{i=p}^m A_i) &= 1- \mathbb{P}(\cap_{i=p}^m A_i^c) = 1- \prod_{i=p}^m \mathbb{P}(A_i^c) \text{ par indépendance}\\
& = 1- \prod_{i=p}^m (1-\mathbb{P}(A_i)) \geq 1- e^{-\sum_{i=p}^m\mathbb{P}(A_i)}
\end{align*} du fait de l'inégalité \(1 - x \leq e^{-x}\) pour
\(x \geq 0\). Ainsi,
\[ \mathbb{P}(\cup_{i=p}^{\infty} A_i) \geq 1- e^{-\sum_{i=p}^{\infty} \mathbb{P}(A_i)} = 1 \]
et l'on conclut finalement que pour tout \(p\),
\(\mathbb{P}(\cup_{i=p}^{\infty} A_i) = 1\), ce qui implique finalement
que \(\mathbb{P}(\lim \sup_{n \to \infty} A_n) = 1\).

\runinhead{Question
3}\phantomsection\addcontentsline{toc}{subsubsection}{Question
3}\label{IV-answer-bc3}
Prendre tous les \(A_n\) égaux à un même événement \(A\) de probabilité
\(\mathbb{P}(A) \in \left]0,1\right[\).

\section{Convergence vers une
constante}\label{IV-convergence-vers-une-constante-1}

\runinhead{Convergence
\(\mathcal{L}^2\)}\phantomsection\addcontentsline{toc}{subsubsection}{Convergence
\(\mathcal{L}^2\)}\label{IV-answer-cvtocst-l2}

\begin{enumerate}
\def\labelenumi{\arabic{enumi}.}
\tightlist
\item
  Supposons \(X_n \overset{\mathcal{L}^2}{\longrightarrow} X\) quand
  \(n\to+\infty\),
  i.e.~\(\mathbb{E}\left(\left(X_n - X\right)^2\right) \xrightarrow[n\to+\infty]{} 0\).
\end{enumerate}

Par positivité de la variance, on a tout d'abord
\[\mathbb{E}\left(\left(X_n - X\right)^2\right) \geq \left(\mathbb{E}\left(X_n-X\right)\right)^2 \geq 0,\]
qui garantit que
\(\left|\mathbb{E}\left(X_n-X\right)\right| = \left|\mathbb{E}(X_n) - \mathbb{E}(X)\right| \to 0\)
quand \(n\to+\infty\).

Ensuite, par inégalité triangulaire on a \begin{align*}
\left|\mathbb{V}\left(X_n\right) - \mathbb{V}\left(X\right)\right| &= \left|\mathbb{E}\left(X_n^2\right) - \mathbb{E}(X_n)^2 - \mathbb{E}\left(X^2\right) + \mathbb{E}(X)^2 \right|\\
& \leq \left|\mathbb{E}\left(X_n^2 - X^2\right)\right| + \left|\mathbb{E}(X_n)^2 - \mathbb{E}(X)^2\right|.
\end{align*}

Or, comme nous avons vu que
\(\mathbb{E}(X_n) \xrightarrow[n\to+\infty]{} \mathbb{E}(X)\), on a
directement que
\(\left|\mathbb{E}(X_n)^2 - \mathbb{E}(X)^2\right| \xrightarrow[n\to+\infty]{} 0\).
Par ailleurs, \begin{align*}
\left|\mathbb{E}\left(X_n^2 - X^2\right)\right| &= \left|\mathbb{E}\left( \left(X_n-X\right)^2 \right) - 2\mathbb{E}\left(X^2\right) + 2\mathbb{E}\left(X\,X_n\right)\right|\\
&= \left|\mathbb{E}\left( \left(X_n-X\right)^2 \right) + 2\mathbb{E}\left(X\,(X_n-X)\right)\right|\\
&\leq \mathbb{E}\left( \left(X_n-X\right)^2 \right) + 2\left|\mathbb{E}\left(X\,(X_n-X)\right)\right|.
\end{align*} Par l'inégalité de Cauchy-Schwartz on a
\[\left|\mathbb{E}\left(X\,(X_n-X)\right)\right|^2 \leq \mathbb{E}\left(X^2\right)\,\mathbb{E}\left( \left(X_n-X\right)^2 \right),\]
donc
\[\left|\mathbb{E}\left(X_n^2 - X^2\right)\right| \leq \mathbb{E}\left( \left(X_n-X\right)^2 \right) + 2\sqrt{\mathbb{E}\left(X^2\right)}\,\sqrt{\mathbb{E}\left( \left(X_n-X\right)^2 \right)},\]
qui tend bien vers \(0\) quand \(n\to+\infty\) par hypothèse. On en
conclut que
\(\left|\mathbb{V}\left(X_n\right) - \mathbb{V}\left(X\right)\right|\xrightarrow[n\to+\infty]{} 0.\)

\begin{enumerate}
\def\labelenumi{\arabic{enumi}.}
\setcounter{enumi}{1}
\tightlist
\item
  Par définition, \(X_n \overset{\mathcal{L}^2}{\longrightarrow} a\)
  quand \(n\to+\infty\) ssi
  \(\mathbb{E}\left(\left(X_n - a\right)^2\right) \xrightarrow[n\to+\infty]{} 0\).
\end{enumerate}

Le premier sens de l'équivalence est immédiatement obtenu par la
question précédente, en prenant \(X = a\) presque-sûrement.

Réciproquement, pour tout \(n\in\mathbb{N}^\ast\), par linéarité de
l'espérance on a \begin{align*}
\mathbb{E}\left( \left(X_n-a\right)^2 \right) &= \mathbb{E}\left( X_n^2 + a^2 - 2aX_n \right) = \mathbb{E}\left(X_n^2\right) + a^2 - 2a\mathbb{E}\left(X_n\right)\\
&= \mathbb{V}\left(X_n\right) + a^2 - 2a\mathbb{E}\left(X_n\right) + \mathbb{E}\left(X_n\right)^2\\
&= \mathbb{V}\left(X_n\right) + \left(\mathbb{E}(X_n) - a\right)^2.
\end{align*} On en déduit immédiatement le deuxième sens de
l'équivalence.

\runinhead{Convergences en loi et
probabilité}\phantomsection\addcontentsline{toc}{subsubsection}{Convergences
en loi et probabilité}\label{IV-answer-cvtocst-loiprob}
Supposons que \(X_n \overset{\mathcal{L}}{\longrightarrow} a\) quand
\(n\to+\infty\). Cela revient à dire que \(X_n\) tend en loi vers la
variable aléatoire \(X = a\) p.s. quand \(n\to+\infty\). Dans ce cas,
\(X\) a pour fonction de répartition
\(x\in\mathbb{R}\mapsto 1_{[a,+\infty[}(x)\) et pour tout
\(x \in \mathbb{R}\backslash\{a\}\), on a
\[\mathbb{P}\left(X_n \leq x \right) \xrightarrow[n\to+\infty]{} 1_{[a,+\infty[}(x).\]

Soit maintenant \(\varepsilon > 0\). D'après la croissance de la
fonction de répartition et l'hypothèse de convergence en loi, on a
\begin{align*}
\mathbb{P}\left(\left|X_n - a\right| \geq \varepsilon \right) &= \mathbb{P}\left(X_n \geq a + \varepsilon \right) + \mathbb{P}\left(X_n \leq a - \varepsilon \right)\\
&= 1 - \mathbb{P}\left(X_n < a + \varepsilon \right) + \mathbb{P}\left(X_n \leq a - \varepsilon \right)\\
&\leq 1 - \mathbb{P}\left(X_n \leq a + \frac{\varepsilon}{2} \right) + \mathbb{P}\left(X_n \leq a - \varepsilon \right)\\
& \xrightarrow[n\to+\infty]{} 1 - 1_{[a,+\infty[}\left(a+\frac{\varepsilon}{2}\right) + 1_{[a,+\infty[}\left(a-\varepsilon\right) = 0.
\end{align*}

Ainsi, on a bien \(X_n \overset{\mathbb{P}}{\longrightarrow} a\) quand
\(n\to+\infty\).

\section{Fonction de répartition
empirique}\label{IV-fonction-de-ruxe9partition-empirique-1}

\runinhead{Préliminaire}\phantomsection\addcontentsline{toc}{subsubsection}{Préliminaire}\label{IV-answer-fdremp-loi}
La variable aléatoire \(1_{]-\infty,x]}(X)\) est à valeurs dans
\(\{0,1\}\) et
\[\mathbb{P}\left( 1_{]-\infty,x]}(X) = 1\right) = F(x) = 1 - \mathbb{P}\left( 1_{]-\infty,x]}(X) = 0\right).\]
On reconnaît une loi de Bernoulli de paramètre \(F(x)\). On en déduit
directement que \(\mathbb{E}(1_{]-\infty,x]}(X)) = F(x)\) et
\(\mathbb{V}(1_{]-\infty,x]}(X)) = F(x)\,(1-F(x))\).

\runinhead{Biais}\phantomsection\addcontentsline{toc}{subsubsection}{Biais}\label{IV-answer-fdremp-biais}
Sous les hypothèses de l'exercice, on a \begin{align*}
\mathbb{E}\left(F_n(x) \right) &= \dfrac{1}{n}\,\sum_{i = 1}^n \mathbb{E}\left( 1_{]-\infty,x]}(X_i) \right) \hspace{1em}\text{par linéarité de l'espérance,}\\
& = \dfrac{1}{n}\,\sum_{i = 1}^n \mathbb{E}\left( 1_{]-\infty,x]}(X) \right) \hspace{1em}\text{car $X_1,\dots,X_n$ ont même loi que $X$,}\\
& = \dfrac{1}{n}\,\sum_{i = 1}^n F(x) \hspace{1em}\text{d'après la première question,}\\
& = F(x).
\end{align*}

\runinhead{Erreur quadratique
moyenne}\phantomsection\addcontentsline{toc}{subsubsection}{Erreur
quadratique moyenne}\label{IV-answer-fdremp-mse}
Soit \(n\in\mathbb{N}^\ast\). D'après la question précédente, on a

\(\mathbb{E}\left( \left(F_n(x) - F(x)\right)^2 \right) = \mathbb{E}\left( \left(F_n(x) - \mathbb{E}(F_n(x))\right)^2 \right) = \mathbb{V}\left(F_n(x)\right).\)

Or \begin{align*}
\mathbb{V}\left(F_n(x)\right) &= \dfrac{1}{n^2}\,\sum_{i = 1}^n \mathbb{V}\left(1_{]-\infty,x]}(X_i)\right) \hspace{1em}\text{par indépendance des variables,}\\
&= \dfrac{1}{n^2}\,\sum_{i = 1}^n \mathbb{V}\left(1_{]-\infty,x]}(X)\right) \hspace{1em}\text{car $X_1,\dots,X_n$ ont même loi que $X$,}\\
&= \dfrac{1}{n^2}\,\sum_{i = 1}^n F(x)\,(1-F(x)) \hspace{1em}\text{d'après la première question,}\\
&= \dfrac{1}{n}\, F(x)\,(1-F(x)) \xrightarrow[n\to+\infty]{} 0.
\end{align*}

On en conclut que
\(F_n(x) \overset{\mathcal{L}^2}{\longrightarrow} F(x)\) quand
\(n\to+\infty\).

\runinhead{Consistance}\phantomsection\addcontentsline{toc}{subsubsection}{Consistance}\label{IV-answer-fdremp-ps}
On remarque que pour tout \(n\in\mathbb{N}^\ast\), \(F_n(x)\) n'est
autre que la moyenne de \(n\) variables aléatoires indépendantes de même
loi de Bernoulli de paramètre \(F(x)\). La loi forte des grands nombres
nous assure donc qu'elle converge presque-sûrement vers l'espérance de
cette loi, qui n'est autre que \(F(x)\).

\addcontentsline{toc}{subsection}{Pour aller plus loin}
\begin{remarque}[{Pour aller plus
loin}]\phantomsection\label{IV-pour-aller-plus-loin}

On peut en fait aller plus loin et montrer que l'on a la convergence
presque sûre uniformément sur \(\mathbb{R}\) voir par exemple
\href{http://math.univ-lyon1.fr/~gelineau/devagreg/Theoreme_Dini.pdf}{ce
document}, c'est le théorème de Glivenko-Cantelli très utile en
statistiques.
\end{remarque}

\section{Théorème de Weierstrass sur
\([0,1]\)}\label{IV-thuxe9oruxe8me-de-weierstrass-sur-01-1}

\runinhead{Preuve du
théorème.}\phantomsection\addcontentsline{toc}{subsubsection}{Preuve du
théorème.}\label{IV-answer-weier-thm}

\begin{enumerate}
\def\labelenumi{\arabic{enumi}.}
\tightlist
\item
  Comme \(X_1,\dots,X_n\) sont indépendantes de même loi de Bernoulli de
  paramètre \(x\in]0,1[\), \(M_n\) est à valeurs dans
  \(\bigl\{\frac{k}{n} : k \in \{0,\dots,n\} \bigr\}\) et pour tout
  \(k\in\{0,\dots,n\}\) on a
  \[\mathbb{P}\left(M_n = \dfrac{k}{n}\right) = \binom{n}{k} x^k\,(1-x)^{n-k}.\]
  Ainsi,
  \[\mathbb{E}\left(f(M_n)\right) = \sum_{k=0}^{n} \binom{n}{k} f\left(\dfrac{k}{n}\right)\, x^k\,(1-x)^{n-k}.\]
\end{enumerate}

Les polynômes de la forme
\(B_{n,k} : u\in [0,1] \mapsto \binom{n}{k} u^k\,(1-u)^{n-k}\) sont
appelés \emph{polynômes de Bernstein}.

\begin{enumerate}
\def\labelenumi{\arabic{enumi}.}
\setcounter{enumi}{1}
\tightlist
\item
  Comme une variable de Bernoulli admet une espérance égale à son
  paramètre, la loi (forte) des grands nombres nous assure que
  \(M_n \to x\) p.s. quand \(n\to+\infty\). Puisque \(f\) est continue,
  on a de même \(f\left(M_n\right) \to f(x)\) p.s. quand
  \(n\to+\infty\).
\end{enumerate}

Par ailleurs, la continuité de \(f\) sur \([0,1]\) nous assure que \(f\)
est bornée (l'image d'un compact par une fonction continue
\(\mathbb{R}\to \mathbb{R}\) est un compact). Ainsi, pour tout
\(n\in\mathbb{N}^\ast\), la variable aléatoire \(f\left(M_n\right)\)
l'est également. Cela nous permet d'appliquer le théorème de convergence
dominée de la question 1~: en notant \((\Omega,\mathcal{A},\mathbb{P})\)
l'espace probabilisé sur lequel sont définies nos variables aléatoires,

\[\mathbb{E}\left( f\left(M_n\right) \right) = \int_{\Omega} f\left(M_n(\omega)\right)\,\mathbb{P}(d\omega) \xrightarrow[n\to+\infty]{} \int_{\Omega} f(x)\,\mathbb{P}(d\omega) = f(x).\]

\begin{enumerate}
\def\labelenumi{\arabic{enumi}.}
\setcounter{enumi}{2}
\item
  En combinant les résultats des deux questions précédentes, on obtient
  que la suite \(\left(P_n\right)_{n\in\mathbb{N}^\ast}\) des polynômes
  définis pour tout \(n\in\mathbb{N}^\ast\) par
  \[P_n : x \in ]0,1[ \mapsto \sum_{k=0}^{n} \binom{n}{k} f\left(\dfrac{k}{n}\right)\, x^k\,(1-x)^{n-k}\]
  converge simplement vers \(f(x)\). On étend simplement ce résultat à
  \([0,1]\) en remarquant que pour tout \(n\in\mathbb{N}^\ast\) on a
  directement l'égalité
  \[P_n(0) = f\left(\dfrac{0}{n}\right) = f(0) \text{ et } P_n(1) = f\left(\dfrac{n}{n}\right) = f(1).\]
  On a même toujours l'égalité
  \(P_n(x) = \mathbb{E}\left( f\left(M_n\right) \right)\) pour
  \(x\in[0,1]\) si l'on remarque qu'une loi de Bernoulli de paramètre
  \(a\in \{0,1\}\) n'est autre d'une Dirac en \(\{a\}\), ce qui signifie
  que \(M_n = a\) p.s. et donc que
  \(\mathbb{E}\left(f(M_n)\right) = f(a) = P_n(a)\).
\item
  On souhaite montrer que
  \[\forall\,\varepsilon > 0 \ \exists\,N_\varepsilon\in\mathbb{N}^\ast : \left(n\geq N_\varepsilon\right) \Rightarrow \left(\forall\,x\in[0,1] : \left|P_n(x) - f(x)\right|\leq \varepsilon\right).\]
\end{enumerate}

Remarquons tout d'abord que la fonction \(f\) étant bornée, elle admet
un maximum sur \([0,1]\), que l'on note \(K\). Puisqu'elle est continue
sur un segment réel, elle est aussi uniformément continue~:
\[\forall\,\varepsilon > 0\ \ \exists\,\delta_{\varepsilon} > 0 : \forall\,(x,y) \in ]0,1[^2,\ \ |x-y| < \delta_\varepsilon \Rightarrow \left|f(x) - f(y)\right| < \varepsilon.\]
Soient \(\varepsilon > 0\), puis \(\delta = \delta_{\varepsilon/2} > 0\)
pour lequel l'implication précédente est vraie avec
\(\dfrac{\varepsilon}{2}\). Alors pour tous \(n\in\mathbb{N}^\ast\) et
\(x\in]0,1[\) on a \begin{align*}
\left|\mathbb{E}\left( f(M_n) \right) - f(x)\right| & \leq \mathbb{E}\left( \left|f(M_n) - f(x) \right| \right) \hspace{1em}\text{par l'inégalité de Jensen,}\\
& = \mathbb{E}\left( \left|f(M_n) - f(x) \right|\,1_{[0,\delta[}\left(\left|M_n - x\right|\right) \right)\\
&\ \ \ + \mathbb{E}\left( \left|f(M_n) - f(x) \right|\,1_{[\delta,+\infty[}\left(\left|M_n - x\right|\right) \right)\\
&\leq \dfrac{\varepsilon}{2} + 2|K|\,\mathbb{P}\left(\left|M_n-x\right|\geq \delta\right) \hspace{1em}\text{par inégalité triangulaire,}\\
&\leq \dfrac{\varepsilon}{2} + 2|K|\,\dfrac{\mathbb{V}(M_n)}{\delta^2} \hspace{1em}\text{par Bienaymé-Chebyshev.}
\end{align*} Or pour tout \(n\in\mathbb{N}^\ast\) \begin{align*}
\mathbb{V}(M_n) &= \mathbb{V}\left(\dfrac{1}{n}\,\sum_{i = 1}^n X_i \right)\\
&= \dfrac{1}{n^2}\,\sum_{i = 1}^n \mathbb{V}(X_i) \hspace{1em}\text{par indépendance de $X_1,\dots,X_n$,}\\
&= \dfrac{1}{n^2}\,\sum_{i = 1}^n x(1-x) \hspace{1em}\text{car $X_1,\dots,X_n$ suivent la même loi de Bernoulli},\\
&= \dfrac{x(1-x)}{n} < \dfrac{1}{4n} \hspace{1em}\text{car $x \in ]0,1[$.}
\end{align*} Par conséquent, pour tout \(n\in\mathbb{N}^\ast\),
\[\left|\mathbb{E}\left( f(M_n) \right) - f(x)\right| < \dfrac{\varepsilon}{2} + \dfrac{|K|}{2n\delta^2}.\]
En prenant \(N_\varepsilon\) le plus petit entier naturel supérieur ou
égal à \(\dfrac{|K|}{\varepsilon\,\delta^2}\), on obtient alors que pour
tout \(n\geq N_\varepsilon\),
\(\left|\mathbb{E}\left( f(M_n) \right) - f(x)\right| < \varepsilon.\)
Comme \(\varepsilon\) et \(N_\varepsilon\) sont les mêmes quel que soit
\(x \in [0,1]\), on obtient bien la convergence uniforme désirée.

\section{Théorème de Slutsky}\label{IV-thuxe9oruxe8me-de-slutsky-1}

\runinhead{Question
1}\phantomsection\addcontentsline{toc}{subsubsection}{Question
1}\label{IV-answer-slut1}
La démonstration est analogue à celle de la
\protect\hyperref[IV-cvceloifdr]{la proposition sur la convergence des
f.d.r. (p.~\pageref*{IV-cvceloifdr})}, où on peut voir que l'on peut
remplacer les fonctions continues bornées \(f_{p,b}\) approchant
\(1_{\left]-\infty,b\right]}\) par des fonctions lipschitziennes
bornées. Voir la démonstration du théorème 18.7 dans Jacod and Protter
(\protect\hyperlink{refIV-Jacod}{2003}) pour une version complète.

\runinhead{Question
2}\phantomsection\addcontentsline{toc}{subsubsection}{Question
2}\label{IV-answer-slut2}
Du fait du résultat de la question 1, il suffit de montrer que
\(\lim_{n \to \infty} \mathbb{E}(f (Y_n)) = \mathbb{E}(f (X))\), pour
toute fonction lipschitzienne bornée. Soit \(f\) une telle fonction. On
a alors \(|f (x)| \leq k\) et \(|f (x) -f (y)| \leq C|x - y|\), pour des
constantes \(k\) et \(C\).

Soit \(\varepsilon > 0\) donné. On a \begin{align*}
|\mathbb{E}(f (Y_n )) - \mathbb{E}(f (X_n ))| &\leq \mathbb{E}(|f (Y_n) - f (X_n )|) \\
                                  &\leq \mathbb{E}(|f (Y_n) - f (X_n )|(1_{|Y_n-X_n|\leq \varepsilon} + 1_{|Y_n-X_n| > \varepsilon}))\\
                                  &\leq C \varepsilon + 2k P(|X_n - Y_n | > \varepsilon)
\end{align*} Le deuxième terme du membre de droite tend vers 0 quand
\(n\) tend vers l'infini, car \((X_n - Y_n )_n\) converge en probabilité
vers 0. Comme \(\varepsilon\) est arbitrairement petit, nous en
déduisons que
\(\lim_{n\to\infty}|\mathbb{E}(f (Y_n )) - \mathbb{E}(f (X_n ))| = 0\),
d'où le résultat.

\runinhead{Question
3}\phantomsection\addcontentsline{toc}{subsubsection}{Question
3}\label{IV-answer-slut3}
Soit \(u \in R^d\) . On a~:
\[|\phi_{Y_n} (u) - \phi_X (u)| \leq |\phi_{Y_n} (u) - \phi_{X_n} (u)| + |\phi_{X_n} (u) - \phi_{X} (u)|\]
D'une part, le \protect\hyperref[IV-levytheorem]{théorème de
Lévy (p.~\pageref*{IV-levytheorem})} partie 1. montre que
\(|\phi_{X_n}(u) - \phi_X (u)|\) tend vers 0 quand \(n \to \infty\).
D'autre part, \begin{align*}
|\phi_{Y_n} (u) - \phi_{X_n} (u)|  & = |\mathbb{E}(e^{i< u, Y_n>} - e^{i< u, X_n>})| \\
& = |\mathbb{E}(e^{i<u,X_n>} (e^{i<u,Y_n-X_n>} - 1))|\leq \mathbb{E}(|e^{i<u,Y_n-X_n>} - 1|)
\end{align*} tend vers 0 quand \(n \to \infty\) d'après
\protect\hyperref[IV-propconv2]{la proposition --- cas
borné (p.~\pageref*{IV-propconv2})} appliquée aux variables aléatoires
\(e^{i<u,Y_n -X_n>}-1\) dont la convergence en proba vers 0 est assurée
par la \protect\hyperref[IV-propconv4]{propriété de
continuité (p.~\pageref*{IV-propconv4})}.
